\documentclass[12pt]{article}
\usepackage[a4paper,margin=1in]{geometry}
\usepackage[T1]{fontenc}
\usepackage{lmodern,microtype}
\usepackage{amsmath,amssymb,amsthm,mathtools}
\usepackage{booktabs,array,enumitem,xcolor}
\usepackage[numbers,sort&compress]{natbib}
\usepackage[colorlinks=true,linkcolor=blue!55!black,citecolor=blue!55!black,urlcolor=blue!55!black]{hyperref}
\linespread{1.07}

\newtheorem{theorem}{Theorem}[section]
\newtheorem{proposition}[theorem]{Proposition}
\theoremstyle{remark}
\newtheorem{remark}[theorem]{Remark}
\newcommand{\norm}[1]{\left\lVert#1\right\rVert}

\title{Physical Edge-Root Matching\\
\large The Length-Four Temporal Clock Resolves a Genuine Spectral Quartet}
\author{Liang Wang}
\date{July 2026}

\begin{document}
\maketitle

\begin{abstract}
RH-192 shows that a rank-four temporal packet cannot be the complete Riesz
shell of the full Frobenius left-multiplication operator.  RH-193 leaves open
a sharper possibility: the packet may resolve four source-excited spectral
modes in the cyclic restriction.  This paper tests that possibility at the
only surviving physical anchor, $\sigma=0.01,L=4$.

The twelve RH-185 windows passing the frozen two-sided $0.10$ residual gate
contain 48 packet roots.  Direct eigendecomposition of the two base physical
operators finds a distinct nearest eigenvalue for every root.  Each
predeclared RH-190 root disc contains exactly one base eigenvalue, the
maximum matching error is below $1.3\times10^{-3}$, and all accepted windows
on a given side select the same four physical modes.  Thus there are eight
unique modes across the left and right channels, not 48 unrelated fits.

For a matched simple eigenvalue with spectral projector $P_\lambda$, the
source and observation states are $P_\lambda S$ and
$P_\lambda^*O^*$.  Every one has nonzero residue
$\operatorname{tr}(OP_\lambda S)$.  The temporal four-space is close to the
canonical span of these source modes: the maximum principal-angle sine is
between $0.035$ and $0.098$ on the right packet and between $0.038$ and
$0.113$ on the left packet.

This is a finite floating-point discovery, not an interval spectral
enclosure or an all-level quartet theorem.  It nevertheless changes the
frontier: the local clock did not manufacture a root grid; it approximated a
genuine source-observable physical edge quartet.
\end{abstract}

\section{Question after the multiplicity correction}

The full matrix state at $\sigma=0.01$ has width $m=256$.  Hence one simple
base eigenvalue produces ambient Frobenius Riesz rank 256, and a literal
rank-one root shell is impossible \citep{WangRH192}.  The source-cyclic
restriction generated by the single matrix $S$, however, can contain one
direction from that repeated eigenspace \citep{WangRH193}.

There are therefore two sharply different explanations of the RH-185 local
candidate:
\begin{enumerate}[leftmargin=2em]
\item the four compressed roots are Ritz artifacts with no nearby physical
 spectrum;
\item they are approximations to four actual base eigenvalues, viewed
 through source-excited matrix directions.
\end{enumerate}
This paper performs the direct finite test.

\section{Frozen data contract}

No window or contour is selected after inspecting the base spectrum.  The
inputs are frozen by RH-185 and RH-190:
\begin{itemize}[leftmargin=2em]
\item scale $\sigma=0.01$;
\item packet length $L=4$;
\item starts 14--18 on the left and 12--18 on the right;
\item two-sided relative residual threshold $0.10$;
\item root-disc radius $0.4\rho\sin(\pi/4)$, where $\rho$ is the source-cycle
 radius of that window.
\end{itemize}
The base dimensions are 512 on the left and 256 on the right.  Both matrix
sources have 256 columns.

For each accepted window we reconstruct the balanced right and left frames
$V,W$, the compressed matrix
\begin{equation}\label{eq:K}
 K=W^*\mathcal L_AV,
\end{equation}
and its four eigenvalues.  We then compute the complete floating spectrum of
the base $A$ and solve a distinct nearest-neighbor assignment.

\section{Root matching and contour count}

Let $\kappa_1,\ldots,\kappa_4$ be the packet roots and let
$\lambda_1,\ldots,\lambda_4$ be their distinct matched base eigenvalues.
The matching error is
\begin{equation}\label{eq:matching-error}
 \epsilon_{\rm match}=\max_j|\kappa_j-\lambda_j|.
\end{equation}

\begin{proposition}[Finite matching result]\label{prop:matching}
Across the twelve accepted windows:
\begin{enumerate}[leftmargin=2em]
\item all 48 packet roots have distinct matched base eigenvalues;
\item every frozen root disc contains exactly one base eigenvalue;
\item $\epsilon_{\rm match}<1.3\times10^{-3}$;
\item all windows on one side select the same four base modes.
\end{enumerate}
\end{proposition}

The final clause is important.  Window motion does not chase different
eigenvalues.  It improves the approximation to one stable edge quartet.

Representative physical eigenvalues have the geometry
\begin{equation}\label{eq:quartet-geometry}
 -0.495\pm0.620i,
 \qquad
 0.474\pm0.557i,
\end{equation}
with small left/right discretization differences.  Their moduli form two
nearby radii, approximately $0.793$ and $0.731$.  This explains why the
compressed roots are close to, but not exactly on, one circle.

\section{Why the result is compatible with RH-192}

For a root disc containing one simple base eigenvalue,
\begin{equation}\label{eq:counts}
 N_A=1,
 \qquad N_{\mathcal L_A}=256.
\end{equation}
The temporal packet count is one.  It agrees with the base count and with a
source-cyclic count, but not with the full Frobenius count.  The finite match
therefore validates the spectral location while confirming the state-type
correction.

Equivalently, if $P_\lambda$ is the rank-one base spectral projector, the
full ambient range is
\begin{equation}\label{eq:ambient-range}
 \{P_\lambda X:X\in\mathbb C^{n\times256}\},
\end{equation}
of dimension 256.  The source orbit selects the single matrix
$P_\lambda S$ from this range.

\section{Source-observation residues}

Let $v_\lambda,w_\lambda$ be right and left eigenvectors normalized by
$w_\lambda^*v_\lambda=1$, so
$P_\lambda=v_\lambda w_\lambda^*$.  Define
\begin{equation}\label{eq:source-observation-states}
 X_\lambda=P_\lambda S,
 \qquad
 Y_\lambda=P_\lambda^*O^*.
\end{equation}
Their Frobenius pairing is
\begin{equation}\label{eq:residue}
 \langle Y_\lambda,X_\lambda\rangle_F
 =\operatorname{tr}(OP_\lambda S)
 =:r_\lambda.
\end{equation}
The same scalar is the residue at $\lambda$ of
\begin{equation}\label{eq:transfer}
 g(z)=\operatorname{tr}\!\left(O(zI-A)^{-1}S\right).
\end{equation}

Every matched mode has $r_\lambda\ne0$ at floating precision.  The minimum
recorded residue modulus across repeated window records is greater than
$1.1\times10^{-2}$.  Therefore all four modes are both source-excited and
observation-visible.  RH-195 proves the exact projector/residue theorem used
here without relying on diagonalizability away from the selected simple
modes.

\section{Canonical spectral subspaces}

For one side, define
\begin{equation}\label{eq:canonical-spaces}
 E_R=\operatorname{span}\{X_{\lambda_1},\ldots,X_{\lambda_4}\},
 \qquad
 E_L=\operatorname{span}\{Y_{\lambda_1},\ldots,Y_{\lambda_4}\}.
\end{equation}
These spaces are exactly invariant under $\mathcal L_A$ and
$\mathcal L_A^*$, respectively.  In contrast, the temporal spaces
$T_R(t),T_L(t)$ are approximate invariant Krylov windows.

For equal-dimensional subspaces $E,F$, write
\begin{equation}\label{eq:gap}
 \operatorname{gap}(E,F)=\sin\theta_{\max}(E,F),
\end{equation}
where $\theta_{\max}$ is the largest principal angle.

\begin{proposition}[Finite subspace alignment]\label{prop:alignment}
For the accepted windows,
\begin{align}
 0.0348&<\operatorname{gap}(T_R,E_R)<0.0978,\\
 0.0384&<\operatorname{gap}(T_L,E_L)<0.1133.
\end{align}
The smallest gaps occur near the latest accepted windows.
\end{proposition}

The alignment is stronger information than eigenvalue proximity alone.  It
shows that the four-dimensional temporal state is approaching the physical
source and observation mode spans, not merely reproducing four scalar roots
through an ill-conditioned compression.

\section{Projector and transversality diagnostics}

The selected base eigenvalues are nonnormal modes.  Their spectral projector
norms are finite but not near one: the maximum is about $17.7$.  The
normalized source-observation overlap
\begin{equation}\label{eq:normalized-overlap}
 \gamma_\lambda=
 \frac{|r_\lambda|}{\norm{X_\lambda}_F\norm{Y_\lambda}_F}
\end{equation}
ranges down to approximately $1.13\times10^{-3}$ on the right channel and
$5.46\times10^{-3}$ on the left.  Thus the canonical packet is genuinely
oblique.  The poor cross-angle conditioning seen in RH-185 is partly
physical, not solely a temporal-window artifact.

This observation also prevents an overstatement.  Root matching does not
automatically provide robust all-level coordinates.  Residues and
transversality must be controlled under refinement.

\section{Determinant and trace comparison}

For the matched quartet define the exact finite spectral polynomial
\begin{equation}\label{eq:physical-determinant}
 d_E(z)=\prod_{j=1}^4(z-\lambda_j).
\end{equation}
The temporal compression gives $d_T(z)=\det(zI-K)$.  We store the determinant
error and the trace-power differences
\begin{equation}\label{eq:trace-errors}
 \operatorname{tr}K^q-\sum_{j=1}^4\lambda_j^q,
 \qquad 1\le q\le8.
\end{equation}
These data quantify more than pointwise root matching and will be used in
RH-199 to formulate a source-channel determinant and trace ledger.

The errors decrease strongly as the temporal subspaces align, but no
asymptotic rate is inferred from seven starts.  The data are a finite
calibration, not a convergence theorem.

\section{Numerical reliability and boundary}

The base eigendecompositions are recomputed once per side using a standard
complex nonsymmetric solver.  Matching is distinct and contour counts are
performed against the complete returned spectrum.  The minimum distance of
any computed eigenvalue from a frozen contour is stored to expose contour
fragility.

What is not supplied is an interval enclosure of the matrix entries,
eigenvalues, or contour winding number.  Therefore ``one eigenvalue per
disc'' is a floating finite statement.  The matching errors are roughly two
orders of magnitude smaller than the contour radii, which makes the signal
clear, but it is not a computer-assisted proof.

Nor is the quartet yet intrinsic.  It is found by comparing a late temporal
window with the complete physical spectrum.  A Gate-A construction must
select and transport the relevant modes without target-dependent matching.

\section{Updated frontier}

The local logical chain is now
\begin{equation*}
 \begin{aligned}
 &\text{temporal packet [approximate, source-generated]}\\
 &\quad\longrightarrow\text{one genuine physical edge quartet [finite]}\\
 &\quad\longrightarrow\text{source-observation spectral channels [next]}\\
 &\quad\longrightarrow\text{canonical packet and transport [open]}.
 \end{aligned}
\end{equation*}
This is a meaningful advance within Gate A.  It does not address the
completion to a unitary/scattering object, a self-adjoint generator, the
$T\log T$ law, von Mangoldt traces, zeta divisors, Hilbert--P\'olya, or RH.

\bibliographystyle{plainnat}
\bibliography{references}
\end{document}
